\documentclass[a4paper,11pt]{article}
\usepackage[a4paper, margin=8em]{geometry}

% usa i pacchetti per la scrittura in italiano
\usepackage[french,italian]{babel}
\usepackage[T1]{fontenc}
\usepackage[utf8]{inputenc}
\frenchspacing 

% usa i pacchetti per la formattazione matematica
\usepackage{amsmath, amssymb, amsthm, amsfonts}

% usa altri pacchetti
\usepackage{gensymb}
\usepackage{hyperref}
\usepackage{standalone}

% imposta il titolo
\title{Appunti Elettronica Digitale}
\author{Luca Seggiani}
\date{2026}

% imposta lo stile
% usa helvetica
\usepackage[scaled]{helvet}
% usa palatino
\usepackage{palatino}
% usa un font monospazio guardabile
\usepackage{lmodern}

\renewcommand{\rmdefault}{ppl}
\renewcommand{\sfdefault}{phv}
\renewcommand{\ttdefault}{lmtt}

% disponi il titolo
\makeatletter
\renewcommand{\maketitle} {
	\begin{center} 
		\begin{minipage}[t]{.8\textwidth}
			\textsf{\huge\bfseries \@title} 
		\end{minipage}%
		\begin{minipage}[t]{.2\textwidth}
			\raggedleft \vspace{-1.65em}
			\textsf{\small \@author} \vfill
			\textsf{\small \@date}
		\end{minipage}
		\par
	\end{center}

	\thispagestyle{empty}
	\pagestyle{fancy}
}
\makeatother

% disponi teoremi
\usepackage{tcolorbox}
\newtcolorbox[auto counter, number within=section]{theorem}[2][]{%
	colback=blue!10, 
	colframe=blue!40!black, 
	sharp corners=northwest,
	fonttitle=\sffamily\bfseries, 
	title=~\thetcbcounter: #2, 
	#1
}

% disponi definizioni
\newtcolorbox[auto counter, number within=section]{definition}[2][]{%
	colback=red!10,
	colframe=red!40!black,
	sharp corners=northwest,
	fonttitle=\sffamily\bfseries,
	title=~\thetcbcounter: #2,
	#1
}

% U.D.M
\newcommand{\amp}{\ensuremath{\mathrm{A}}}
\newcommand{\volt}{\ensuremath{\mathrm{V}}}
\newcommand{\meter}{\ensuremath{\mathrm{m}}}
\newcommand{\second}{\ensuremath{\mathrm{s}}}
\newcommand{\farad}{\ensuremath{\mathrm{F}}}
\newcommand{\henry}{\ensuremath{\mathrm{H}}}
\newcommand{\siemens}{\ensuremath{\mathrm{S}}}

% circuiti
\usepackage{circuitikz}
\usetikzlibrary{babel}

% disegni
\usepackage{pgfplots}
\pgfplotsset{width=10cm,compat=1.9}

% disponi codice
\usepackage{listings}
\usepackage[table]{xcolor}

\definecolor{codegreen}{rgb}{0,0.6,0}
\definecolor{codegray}{rgb}{0.5,0.5,0.5}
\definecolor{codepurple}{rgb}{0.58,0,0.82}
\definecolor{backcolour}{rgb}{0.95,0.95,0.92}

\lstdefinestyle{codestyle}{
		backgroundcolor=\color{black!5}, 
		commentstyle=\color{codegreen},
		keywordstyle=\bfseries\color{magenta},
		numberstyle=\sffamily\tiny\color{black!60},
		stringstyle=\color{green!50!black},
		basicstyle=\ttfamily\footnotesize,
		breakatwhitespace=false,         
		breaklines=true,                 
		captionpos=b,                    
		keepspaces=true,                 
		numbers=left,                    
		numbersep=5pt,                  
		showspaces=false,                
		showstringspaces=false,
		showtabs=false,                  
		tabsize=2
}

\lstdefinestyle{shellstyle}{
		backgroundcolor=\color{black!5}, 
		basicstyle=\ttfamily\footnotesize\color{black}, 
		commentstyle=\color{black}, 
		keywordstyle=\color{black},
		numberstyle=\color{black!5},
		stringstyle=\color{black}, 
		showspaces=false,
		showstringspaces=false, 
		showtabs=false, 
		tabsize=2, 
		numbers=none, 
		breaklines=true
}

\lstdefinelanguage{javascript}{
	keywords={typeof, new, true, false, catch, function, return, null, catch, switch, var, if, in, while, do, else, case, break},
	keywordstyle=\color{blue}\bfseries,
	ndkeywords={class, export, boolean, throw, implements, import, this},
	ndkeywordstyle=\color{darkgray}\bfseries,
	identifierstyle=\color{black},
	sensitive=false,
	comment=[l]{//},
	morecomment=[s]{/*}{*/},
	commentstyle=\color{purple}\ttfamily,
	stringstyle=\color{red}\ttfamily,
	morestring=[b]',
	morestring=[b]"
}

% disponi sezioni
\usepackage{titlesec}

\titleformat{\section}
	{\sffamily\Large\bfseries} 
	{\thesection}{1em}{} 
\titleformat{\subsection}
	{\sffamily\large\bfseries}   
	{\thesubsection}{1em}{} 
\titleformat{\subsubsection}
	{\sffamily\normalsize\bfseries} 
	{\thesubsubsection}{1em}{}

% disponi alberi
\usepackage{forest}

\forestset{
	rectstyle/.style={
		for tree={rectangle,draw,font=\large\sffamily}
	},
	roundstyle/.style={
		for tree={circle,draw,font=\large}
	}
}

% disponi algoritmi
\usepackage{algorithm}
\usepackage{algorithmic}
\makeatletter
\renewcommand{\ALG@name}{Algoritmo}
\makeatother

% disponi numeri di pagina
\usepackage{fancyhdr}
\fancyhf{} 
\fancyfoot[L]{\sffamily{\thepage}}

\makeatletter
\fancyhead[L]{\raisebox{1ex}[0pt][0pt]{\sffamily{\@title \ \@date}}} 
\fancyhead[R]{\raisebox{1ex}[0pt][0pt]{\sffamily{\@author}}}
\makeatother

\begin{document}
% sezione (data)
\section{Lezione del 27-04-26}

% stili pagina
\thispagestyle{empty}
\pagestyle{fancy}

% testo
Nella scorsa lezione abbiamo visto l'amplificatore non invertente basato sull'amplificatore operazionale. Vediamo adesso la configurazione invertente.

\subsubsection{Amplificatore invertente}
Vediamo quindi la configurazione \textit{invertente} dell'amplificatore:

\begin{center}
\begin{tikzpicture}[transform shape]
	% Paths, nodes and wires:
	\node[op amp] at (0.19, 0.51){};
	\draw (-1, 0.02) -| (-2, -0) -- (-2, -1);
	\draw (-1, 1) to[american resistor, l_={$R_1$}] (-3, 1);
	\draw (-3, 1) to[american voltage source, l_={$v_s$}] (-3, -1);
	\node[tlground] at (-3, -1){};
	\node[tlground] at (-2, -1){};
	\draw (-1, 1) -| (-1, 3);
	\draw (-1, 3) to[american resistor, l_={$R_2$}] (1.25, 3);
	\draw (1.25, 3) -| (1.38, 0.51);
	\draw (1.38, 0.51) -| (3, 0.5);
	\node[shape=rectangle, minimum width=0.465cm, minimum height=0.465cm] at (3.25, 0.5){} node[anchor=center, align=center, text width=0.077cm, inner sep=6pt] at (3.25, 0.5){$v_u$};
\end{tikzpicture}
\end{center}

In questo caso, come vediamo, basta applicare il segnale all'entrata invertente ($-$) anziché non invertente ($+$) dell'OP amp.

Possiamo quindi svolgere l'analisi come nel caso invertente, prendendo l'approssimazione a cortocircuito virtuale:
$$
\begin{cases}
v^+ \approx v^- \\
i^+ \approx i^- \approx 0
\end{cases}
$$

Sulla $R_1$, presa la direzione della corrente da sinistra verso destra, avremo:
$$
i_1 = \frac{V_s}{R_1}
$$
mentre sulla $R_2$ aremo:
$$
i_2 = i_1, \quad v_o = - R_2 i _ = - R_2 i_1 = -R_2 \frac{V_s}{R_1}
$$
per cui:
$$
V_o = -\frac{R_2}{R_1} V_s
$$
Come nel caso invertente, questo dipende solo dalla rete di condizionamento esterna, e non dalle caratteristiche intrinseche all'OP amp, né da altre caratteristiche variabili come la temperatura. 

Per quanto riguarda l'impedenza di uscita, usiamo il risultato dalla teoria della reazione:
$$
R_{of} = \frac{R_{OPA}}{1 - \beta A} \approx 0
$$
mentre per l'impedenza di ingresso prendiamo la sola resistenza $R_1$, che è quella vista dal segnale:
$$
R_i = R_1
$$

\subsubsection{Amplificatore di differenza}
Dopo aver visto le configurazioni invertenti e non dell'amplificatore operazionale, vediamo come si può impiegare questo componente per amplificare la \textit{differenza} fra due segnali.

\begin{center}
\begin{tikzpicture}[transform shape]
	% Paths, nodes and wires:
	\node[op amp] at (0.19, 0.51){};
	\draw (-1, 1) to[american resistor, l_={$R_1$}] (-3, 1);
	\draw (-5, 1) to[american voltage source, l_={$v_2$}] (-5, -1);
	\node[tlground] at (-5, -1){};
	\draw (-1, 1) -| (-1, 3);
	\draw (-1, 3) to[american resistor, l_={$R_2$}] (1.25, 3);
	\draw (1.25, 3) -| (1.38, 0.51);
	\draw (1.38, 0.51) -| (3, 0.5);
	\node[shape=rectangle, minimum width=0.465cm, minimum height=0.465cm] at (3.25, 0.5){} node[anchor=center, align=center, text width=0.077cm, inner sep=6pt] at (3.25, 0.5){$v_u$};
	\draw (-5, 1) -- (-3, 1);
	\draw (-1, 0.02) to[american resistor, l={$R_3$}] (-3, -0);
	\draw (-1, -0) to[american resistor, l={$R_4$}] (-1, -2);
	\draw (-3, -0) to[american voltage source, l_={$v_1$}] (-3, -2);
	\node[tlground] at (-3, -2){};
	\node[tlground] at (-1, -2){};
\end{tikzpicture}
\end{center}

Questa configurazione è equivalente a quella con entrata invertente, ma con un secondo segnale ($v_1$) che viene ammesso, dopo un partitore di tensione, all'ingresso non invertente dell'amplificatore.

Per studiare questo circuito, prendiamo più casi separati:
\begin{itemize}
\item $V_1 \neq 0$, $V_2 = 0$ \\
In questo caso la $V_1$ va a ground, per cui se non per il partitore $R_3$ $R_3$ siamo nella configurazione non invertente. Quindi:
$$
V^+ = V_1 \frac{R_4}{R_3 + R_4}
$$
e semplicemente applicando la legge di uscita dell'amplificatore non invertente abbiamo:
$$
V_{o1} = V^* \left(1 + \frac{R_2}{R_1}\right) = V_1 \frac{R_4}{R_3 + R_4} \left(1 + \frac{R_2}{R_1}\right)
$$

\item $V_1 = 0$, $V_2 \neq 0$ \\

In questo è $V_2$ ad andare a ground, per la tensione all'ingresso non invertente sarà nulla:
$$
V^+ = 0
$$
per cui ci troveremo praticamente nella configurazione invertente, dove invece di una resistenza avevamo il ramo non invertente direttamente a ground. Possimao quindi applicare la legge di uscita dell'amplificatore invertente per avere:
$$
V_{o2} = -\frac{R_2}{R_1} V_2
$$

\end{itemize}

Sommiamo i contributi trovati per trovare la tensione di uscita complessiva del circuito (sovrapposizione degli effetti):
$$
V_o = V_{o1} + V_{o2} = V_1 \frac{R_4}{R_3 + R_4} (1 + \frac{R_2}{R_1}) - V_2 \frac{R_2}{R_1}
$$

Notiamo che questo è ancora distante da ciò che vorremmo, ovvero:
$$
V_o = K (V_1 - V_2)
$$
Cioè vogliamo bilanciare le resistenze perché i contributi di $V_1$ e $V_2$ si eguaglino. Ciò che facciamo allora è imporre:
$$
V_1 = V_2 \implies V_o = 0
$$
che diventa:
$$
\frac{R_4}{R_3 + R_4} \left( 1 + \frac{R_2}{R_1} \right) = \frac{R_2}{R_1} \implies \frac{R_4}{R_3 + R_4} \frac{R_1 + R_2}{R_1} = \frac{R_2}{R_1} 
$$
$$
\implies \frac{R_4}{R_3 + R_4} = \frac{R_2}{R_1 + R_2} \implies \frac{R_3 + R_4}{R_4} = \frac{R_1 + R_2}{R_2} \implies \frac{R_3}{R_4} = \frac{R_1}{R_2}
$$
cioè c'è un intera famiglia di soluzioni dell'equazione che ci permettono di ottenere amplificatori di differenza. Resterà da calcolare il guadagno $K$ sulla differenza a partire dalla soluzione scelta. Questo si fa sostituendo la soluzione trovata:
$$
\frac{R_4}{R_3 + R_4} = \frac{R_2}{R_1 + R_2}
$$
nell'equazione della $V_o$:
$$
V_o = V_1 \frac{R_2}{R_1 + R_2} \left( \frac{R_1 + R_2}{R_1} \right) - V_2 \frac{R_2}{R_1} = \frac{R_2}{R_1} \left( V_1 - V_2 \right) \implies K = \frac{R_2}{R_1}
$$
cioè il guadagno è sempre $\frac{R_2}{R_1}$, indipendente dall'OP amp e deciso dalla rete esterna.

Veniamo quindi alla resistenza in ingresso. Questa si calcolerà separatemente per $V_1$ e $V_2$:
\begin{itemize}
\item Per $V_1$, si ha che $R_{i1}$ è:
$$
R_{i1} = \frac{v_1}{i_1} \Big|_{V_2 = 0} = R_3 + R_4
$$
in quanto $i_1$ vale:
$$
i_1 = \frac{V_1}{R_3 + R_4}
$$
\item Per $V_2$, si ha che $R_{i2}$ è:
$$
R_{i2} = \frac{v_2}{i_2} \Big|_{V_1 = 0} = R_1
$$
in quanto $i_2$ vale:
$$
i_2 = \frac{v_2}{R_1}
$$
\end{itemize}

$R_o$ tende invece a zero per le caratteristiche dell'OP amp e la teoria della reazione: 
$$
R_o \rightarrow 0
$$

\subsubsection{Amplificatore di somma}
Vediamo il circuito \textit{sommatore}, che ha il compito di (appunto) sommare due segnali attraverso un amplificatore operazionale.

\begin{center}
\begin{tikzpicture}[transform shape]
	% Paths, nodes and wires:
	\node[op amp] at (0.19, 0.51){};
	\draw (-2, 1) to[american resistor, l_={$R_1$}] (-4, 1);
	\draw (-6, 1) to[american voltage source, l_={$v_2$}] (-6, -1);
	\node[tlground] at (-6, -1){};
	\draw (-1, 1) -| (-1, 3);
	\draw (-1, 3) to[american resistor, l_={$R_2$}] (1.25, 3);
	\draw (1.25, 3) -| (1.38, 0.51);
	\draw (1.38, 0.51) -| (3, 0.5);
	\node[shape=rectangle, minimum width=0.465cm, minimum height=0.465cm] at (3.25, 0.5){} node[anchor=center, align=center, text width=0.077cm, inner sep=6pt] at (3.25, 0.5){$v_u$};
	\draw (-4, 1) to[american voltage source, l_={$v_1$}] (-4, -1);
	\node[tlground] at (-4, -1){};
	\node[tlground] at (-2, -1){};
	\draw (-1, -0) -- (-2, -0) -| (-2, -1);
	\draw (-2, 1) -- (-1, 1);
	\draw (-2, 3) to[american resistor, l={$R_3$}] (-4, 3);
	\draw (-2, 3) -- (-2, 1);
	\draw (-4, 3) -| (-6, 1);
\end{tikzpicture}
\end{center}

Vediamo come in questo caso introduciamo entrambi i segnali all'ingresso non invertente dell'OP amp. Le correnti entranti saranno quindi:
$$
i_1 = \frac{V_1}{R_1}, \quad i_3 = \frac{V_2}{R_3} \implies i_1 + i_3 = i_2 + i^-
$$
con $i_2$ presa diretta da sinistra verso destra. Con l'approssimazione del cortocircuito virtuale si ha quindi:
$$
i_2 = i_1 + i_3
$$
A questo punto l'uscita è data dalla caduta su $R_2$:
$$
V_o = -R_2 i_2 = -R_2 \left( \frac{V_1}{R_1} + \frac{V_2}{R_3} \right)
$$
e notiamo che basterà scegliere $R_1 = R_3 = R$ per avere un amplificazione:
$$
V_o = -\frac{R_2}{R} (V_1 + V_2)
$$
invertente della somma dei due segnali.

Notiamo che il cortocircuito virtuale mantiene il polo condiviso da $R_3$ e $R_1$ a 0, per cui non abbiamo problematiche di interazione fra i generatori $v_1$ e $v_2$.

\subsubsection{Integratore di Miller}
Attraverso gli amplificatori operazionali si possono realizzare \textit{integratori} analogici, cioè circuiti che calcolano gli integrali di segnali in entrata.

\begin{center}
\begin{tikzpicture}[transform shape]
	% Paths, nodes and wires:
	\draw (-2, 1) to[american resistor, l_={$R$}] (-4, 1);
	\draw (-1, 1) -| (-1, 3);
	\draw (1.25, 3) -| (1.38, 0.51);
	\draw (1.38, 0.51) -| (3, 0.5);
	\node[shape=rectangle, minimum width=0.465cm, minimum height=0.465cm] at (3.25, 0.5){} node[anchor=center, align=center, text width=0.077cm, inner sep=6pt] at (3.25, 0.5){$v_u$};
	\draw (-4, 1) to[american voltage source, l_={$v_s$}] (-4, -1);
	\node[tlground] at (-4, -1){};
	\node[tlground] at (-2, -1){};
	\draw (-1, -0) -- (-2, -0) -| (-2, -1);
	\draw (-2, 1) -- (-1, 1);
	\draw (-1, 3) to[capacitor, l={$C$}] (1.25, 3);
	\node[op amp] at (0.19, 0.51){};
\end{tikzpicture}
\end{center}

Per valutare se questo circuito va effettivamente a comportarsi come integratore, valutiamo la funzione di trasferimento:
$$
H = \frac{V_o}{V_s}
$$

Abbiamo 2 modalità per fare ciò: nel dominio di Laplace e nel dominio tempo. Iniziamo con l'analisi nel dominio di Laplace:
\begin{itemize}
\item \textbf{\textsf{Dominio di Laplace}} \\
	Nel dominio di Laplace avremo il generatore di segnale in funzione del parametro $s$, cioè $V_s(s)$, la resistenza resterà $R$ e la capacità sarà espressa come $\frac{1}{Cs}$. La configurazione dell'OP amp è quella ad amplificatore invertente, perciò avremo la funzione per l'uscita $V_o(s)$:
$$
V_o(s) = - \frac{\frac{1}{Cs}}{R} V_s(s) = - \frac{1}{RC s} V_s(s)
$$
da cui in dominio tempo:
$$
V_o(t) = -\frac{1}{RC} \int_0^t V_s(\tau) \, d\tau + V_o(o)
$$
Da questo segue che la funzione di trasferimento sarà:
$$
H(s) = -\frac{1}{RCs}
$$
ovvero l'integratore moltiplicato per $\frac{1}{RC}$.

\item \textbf{\textsf{Dominio tempo}} \\
Per fare la stessa cosa in dominio tempo applichiamo la caratteristica in dominio tempo del condensatore, noto il fatto che la corrente su $R$ è:
$$
i_R = \frac{V_s}{R} = i_C
$$
e che per l'approssimazione a cortocircuito virtuale la stessa corrente passa su $i_C$. Quindi si avrà:
$$
V_C(t) = \frac{1}{C} \int_0^t i_C(t) \, dt + V_C(0) = \frac{1}{RC} \int_0^t V_s(t) \, dt + V_C(0)
$$
e visto che $V_o = - V_C$:
$$
V_o(t) = - \frac{1}{RC} \int_0^t V_s(t) \, dt + V_o(0)
$$
che è la stessa caratteristica di prima.
\end{itemize}

Notiamo che questo circuito, avendo un polo all'origine, non è stabile. Di base un'integratore non è mai stabile, in quanto ad esempio integrando una costante si ottiene una rampa che teoricamente arriva ad infinito.
Notiamo che nel caso reale prima o poi l'OP amp arriverà a saturazione, cioè la tensione raggiungerà una delle due rotaie e di li in poi sarà costretta a non superare un valore costante. Si può risolvere il problema rendendo l'integratore \textit{leaky} mettendo una resistenza in serie alla capacità, quindi limitandone la banda ma aumentandone la stabilità.

Inoltre, scambiando capacità e resistenza si può trasformare il circuito in un \textit{derivatore}, che però non vediamo in quanto è un oggetto ancora più instabile dell'integratore. 

\subsubsection{Regolatore di tensione lineare serie}
Vediamo un approccio alla \textit{regolazione di tensione} che supera lo Zener (solitamente incapace di fornire grandi potenze). Data una sorgente con offset DC e una variazione AC sovrapposta, vogliamo inserire un \textit{elemento di passo} col compito di rimuovere corrente quando c'è un surplus, e fornirne quando ci suono dei vuoti.
\begin{center}
\begin{tikzpicture}[transform shape]
	% Paths, nodes and wires:
	\draw (-5, -0) to[battery1, l_={$E$}] (-5, -2);
	\draw (-5, 2) to[american voltage source, l_={$V_s(t)$}] (-5, -0);
	\node[tlground] at (-5, -2){};
	\draw (-5, 2) -| (-5, 3) -- (-3, 3);
	\node[shape=rectangle, draw, line width=1pt, minimum width=3.965cm, minimum height=1.965cm] at (-1, 3){} node[anchor=center, align=center, text width=3.577cm, inner sep=6pt] at (-1, 3){Elemento di passo};
	\draw (1, 3) -- (3, 3) -| (3, 2);
	\draw (3, 2) to[american resistor, l={$R_L$}] (3, -2);
	\node[tlground] at (3, -2){};
	\node[tlground] at (2.957, -2.062){};
\end{tikzpicture}
\end{center}

Notiamo che l'elemento di passo è un \textit{elemento di potenza}, cioè un elemento progettato per fornire potenze non trascurabili. Solitamente, questi vengono realizzati attraverso transistori bipolari o MOSFET appositamente progettati. Noi prenderemo l'esempio del bipolare. 

Infine, per poter utilizzare l'elemento di passo, vorremo confrontare la tensione in uscita con un certo riferimento, attraverso un \textit{comparatore} (scopo per cui l'amplificatore operazionale che abbiamo studiato finora è perfetto):
\begin{center}
\begin{tikzpicture}[transform shape]
	% Paths, nodes and wires:
	\draw (-5, -0) to[battery1, l_={$E$}] (-5, -2);
	\draw (-5, 2) to[american voltage source, l_={$V_s(t)$}] (-5, -0);
	\node[tlground] at (-5, -2){};
	\draw (-5, 2) -| (-5, 3) -- (-3, 3);
	\node[shape=rectangle, draw, line width=1pt, minimum width=3.965cm, minimum height=1.965cm] at (-1, 3){} node[anchor=center, align=center, text width=3.577cm, inner sep=6pt] at (-1, 3){Elemento di passo};
	\draw (1, 3) -- (3, 3) -| (3, 2);
	\draw (3, 2) to[american resistor, l={$R_L$}] (3, -2);
	\node[tlground] at (3, -2){};
	\node[tlground] at (2.957, -2.062){};
	\node[shape=rectangle, draw, line width=1pt, minimum width=1.965cm, minimum height=0.965cm] at (0.5, -0.5){} node[anchor=center, align=center, text width=1.577cm, inner sep=6pt] at (0.5, -0.5){\scriptsize Comparatore};
	\draw[-latex] (-0.5, -0.5) -| (-1, 2);
	\draw[dash pattern={on 1.6pt off 1.6pt}, -latex] (2.5, 3) |- (1.5, -0.25);
	\draw (2, -1.75) |- (1.5, -0.75);
	\node[shape=rectangle, minimum width=0.965cm, minimum height=0.465cm] at (2, -2){} node[anchor=north west, align=left, text width=0.577cm, inner sep=6pt] at (1.5, -1.75){$V_{ref}$};
\end{tikzpicture}
\end{center}

\newpage

Il circuito che consideriamo è quindi il seguente:
\begin{center}
\begin{tikzpicture}[transform shape]
	% Paths, nodes and wires:
	\draw (-5, -0) to[battery1, l_={$E$}] (-5, -2);
	\draw (-5, 2) to[american voltage source, l_={$V_s(t)$}] (-5, -0);
	\node[tlground] at (-5, -2){};
	\draw (-5, 2) -| (-5, 3) -- (-3, 3);
	\draw (1, 3) -- (3, 3) -| (3, 2);
	\draw (3, 2) to[american resistor, l={$R_L$}] (3, -2);
	\node[tlground] at (3, -2){};
	\node[tlground] at (2, -2){};
	\node[npn, rotate=90] at (-0.98, 3){};
	\draw (-3, 3) -- (-1.75, 3);
	\draw (-0.21, 3) -- (1, 3);
	\node[op amp, xscale=-1] at (0.31, 0.01){};
	\draw (1.5, 0.5) -- (2, 0.5);
	\draw (2, 0.5) to[american resistor, l_={$R_1$}] (2, 3);
	\draw (2, 0.5) to[american resistor, l={$R_2$}] (2, -2);
	\draw (-0.88, 0.01) -| (-1, -0) -| (-0.98, 2.16);
	\draw (-4, -2) to[empty Zener diode] (-4, -0);
	\draw (-4, -0) to[american resistor, l_={$R_Z$}] (-4, 3);
	\node[tlground] at (-4, -2){};
	\draw (1.5, -0.5) -| (1.5, -1.5) -- (-2, -1.5) -| (-2, -0) |- (-4, -0);
	\node[shape=rectangle, minimum width=0.465cm, minimum height=0.465cm] at (-1.5, 3.25){} node[anchor=center, align=center, text width=0.077cm, inner sep=6pt] at (-1.5, 3.25){C};
	\node[shape=rectangle, minimum width=0.465cm, minimum height=0.465cm] at (-0.46, 3.25){} node[anchor=center, align=center, text width=0.077cm, inner sep=6pt] at (-0.46, 3.25){E};
\end{tikzpicture}
\end{center}

Dove l'idea è quella di modulare un transistore BJT con la variazione rispetto ad un certo voltaggio di riferimento, calcolato attraverso il diodo di Zener a sinistra. A destra (cioè all'ingresso invertente) avremo un campionamento della tensione di uscita, appositamente partizionata.

Le ipotesi che facciamo sono quindi quelle di:
\begin{itemize}
	\item Cortocircuito virtuale per l'OP amp, che si trova in reazione negativa e in regime lineare;
	\item Diodo di Zener in polarizzazione negativa.
\end{itemize}

Possiamo quindi prendere la $V^-$ e la $V^+$:
$$
\begin{cases}
	V^- = V_o \frac{R_2}{R_1 + R_2} \\
	V^+ = V_z
\end{cases}
$$
da cui immediatamente si ha che:
$$
V_o \frac{R_2}{R_1 + R_2} = V_z \implies V_o = V_z \frac{R_1 + R_2}{R_2}
$$
cioè siamo riusciti ad ottenere un componente che fornisce una tensione di uscita condizionata da $R_1$ e $R_2$, e dipendente dal valore di riferimento fornito dallo Zener.

Vediamo cosa accade in tempo reale durante l'operazione dell'elemento di passo.
\begin{enumerate}
	\item Se $V_o$ aumenta, di conseguenza $V^-$ aumenta (in quanto ne dipende direttamente, attraverso il partitore). Col fatto che $V^+$ resta costante (dallo Zener), la $V_u = K(V^+ - V^-)$ in uscita dall'operazionale diminuisce, e quindi abbassiamo la corrente alla base del BJT. Conseguenza diretta è che la $V_o$ viene portata a diminuire, e quindi la reazione stabilizza la variazione di tensione.
\end{enumerate}

\end{document}

